% LSPV-5 REPORT::AFFINE-TO-QUADRATIC-SMALLBIAS-FIREWALL237::2026-10-31T01:48:00-07:00
\documentclass[11pt]{article}
\usepackage[margin=.86in]{geometry}
\usepackage{amsmath,amssymb,amsthm,booktabs,xurl,hyperref}
\newtheorem{theorem}{Theorem}
\newtheorem{corollary}{Corollary}
\title{LSPV-5 AFFINE-TO-QUADRATIC-SMALLBIAS-FIREWALL237\\
An Exact Bent-Quadratic Countercontrol and Joint ARC Referee Ledger}
\author{LSPV-5 Independent Red Team / Synthesis}
\date{31 October 2026, 01:48 America/Los\_Angeles (PDT)}
\begin{document}\maketitle

\noindent\textbf{Status.}
\textbf{PROVED/EXACT} for the bent-quadratic small-bias theorem below;
\textbf{COMPUTATIONAL} for finite Fourier replay and matched original-source
SAT/UNSAT controls; \textbf{FALSIFIED/RETIRED} is any extension of the
AFFINE-SMALLBIAS proof to quadratic maps using only linear small bias.
The archived affine theorem itself remains correct under its stated
source-independent affine hypotheses.  Universal arbitrary-$A\mapsto\Phi(A)$,
ARC-EDGE-90, $P=NP$, and $P\ne NP$ remain \textbf{OPEN}.

\section{Target, quantifiers, and four falsifiers}
The target was not a new SAT route.  It was the exact scope of the load-bearing
AFFINE-SMALLBIAS-DESCENT-GATE2 argument.  Its Fourier proof controls inverse
images of affine output sets because affine fibers have indicators that are
linear-character sums.  The missing question was whether the same
``small-bias therefore no descent'' inference survives even the first nonlinear
degree.

The claim card is: for every $r\ge1$, construct one explicit distribution on
$2r$ bits which has vanishing linear bias but is separated from uniform by one
fixed source-independent quadratic bit.  The model is exact finite probability
over $\mathbb F_2$; no SAT oracle, circuit lower bound, or unknown guest is used.
The SAME\_REALIZATION objects are the same random vector for every linear test
and for the quadratic test.

Four falsifiers were posed.  F1: linear $\epsilon$-bias controls every quadratic
output bit---false by the theorem.  F2: this refutes the archived affine
theorem---false, because that theorem never claims nonlinear control.  F3: the
constructed distribution is automatically the CHMY easy-source ensemble---
false; no such identification is made.  F4: a nonlinear semantic escape is a
fully charged ARC constructor---false; label preservation, acquisition,
serialization, access, and native-machine resources remain separate obligations.

\section{Exact quadratic firewall}
Write $x=(u_1,v_1,\ldots,u_r,v_r)\in\mathbb F_2^{2r}$ and define
\[
 q_r(x)=\bigoplus_{i=1}^r u_i v_i.
\]
Let $D_r$ be uniform on $S_r=\{x:q_r(x)=0\}$, and let $U_r$ be uniform on all
of $\mathbb F_2^{2r}$.

\begin{theorem}[Linear-small-bias does not fool one quadratic bit]
For every nonzero $a\in\mathbb F_2^{2r}$,
\[
 \left|\mathbb E_{D_r}(-1)^{a\cdot X}\right|=\frac{1}{2^r+1}.
\]
Thus $D_r$ is $\epsilon_r$-biased with $\epsilon_r=(2^r+1)^{-1}\to0$.
Nevertheless the fixed quadratic map $R_r(x)=q_r(x)$ satisfies
\[
 \Pr_{D_r}[R_r=1]=0,
 \qquad
 \Pr_{U_r}[R_r=1]=\frac{1-2^{-r}}{2},
\]
so its distinguishing gap tends to $1/2$.
\end{theorem}
\begin{proof}
For one pair,
\[
 \sum_{u,v\in\mathbb F_2}(-1)^{uv+\alpha u+\beta v}
 =2(-1)^{\alpha\beta}.
\]
Multiplication over the $r$ independent pairs gives the unnormalised Walsh
coefficient $\sum_x(-1)^{q_r(x)+a\cdot x}=\pm2^r$.  At $a=0$ this yields
\[
 |S_r|=2^{2r-1}+2^{r-1},\qquad
 |q_r^{-1}(1)|=2^{2r-1}-2^{r-1}.
\]
For $a\ne0$, $\sum_x(-1)^{a\cdot x}=0$, while
$1_{S_r}(x)=(1+(-1)^{q_r(x)})/2$.  Hence
\[
 \left|\sum_{x\in S_r}(-1)^{a\cdot x}\right|=2^{r-1},
\]
and division by $|S_r|$ gives $(2^r+1)^{-1}$.  The second display follows
from the two exact level-set sizes.
\end{proof}

\begin{corollary}[Exact proof-method boundary]
No argument whose sole pseudorandomness premise is small bias against nonzero
linear characters can bound all source-independent quadratic output maps.  A
quadratic extension needs an independently proved quadratic-fooling property
(or stronger), plus the original label and construction obligations.
\end{corollary}

This is the familiar bent-quadratic phenomenon, not a novelty claim.  Rothaus
introduced bent functions and their flat Walsh spectra.  Naor--Naor small-bias
spaces control parities; that definition does not promise fooling quadratic
tests.

\section{Referee audit of AFFINE-SMALLBIAS-DESCENT-GATE2}
The affine report's corrected counting statement remains sound.  If the target
truth table has $m$ input bits and target circuits have size at most $u_N$, the
number $B_N$ of target YES descriptions obeys
\[
 \log_2 B_N=O((u_N+1)\log(m+u_N+2))\le N^{\nu+o(1)}.
\]
With its cited easy-source distribution of bias
$\epsilon_N=2^{-\lceil N^\rho\rceil}$ for $\nu<\rho<\mu$, one has
$B_N\epsilon_N=o(1)$.  For each fixed affine map, Fourier expansion of an affine
fiber compares the image of that easy distribution with the image of uniform
source bits.  The proof does \emph{not} replace the latter image by a uniform
target distribution.  Averaging the pointwise fixed-seed inequality over a seed
independent of the source is valid.  The YES/NO density and set-size quantifiers
in the archived handoff are therefore preserved.

What fails is promotion beyond that exact scope.  The theorem above shows that
an arbitrarily small linear bias by itself gives no useful bound on even one
quadratic output coordinate.  It does not show that the particular CHMY local
generator is broken by $q_r$, that its support equals $S_r$, or that a quadratic
MCSP reduction exists.  It only identifies the additional pseudorandomness
lemma that a nonlinear extension would require.

\section{Joint ARC/FLMV reconciliation}
LSPV-4's newest recurrence algebra is correct conditional on its unconstructed
native-machine hypotheses.  Writing child-count, guest-emulation, and length
perturbations as $\theta,\beta,\xi$, the strict time margin at the FLMV test point
is
\[
 \theta+\beta+{241\over200}\xi<{1\over400},
\]
and the physical child-space margin is
\[
 \xi<{26279\over23232400}.
\]
These simultaneous inequalities cannot be populated by the affine theorem or
by the quadratic countercontrol.  They still require one source-computable,
both-label, hereditary ordinary-SAT constructor with all selector discovery,
failed policies, source accesses, physical output or virtual access, reset,
guest execution, proof checking, and ancestor state charged on the actual FLMV
indexed machine.

Current specialist results remain restricted: LSPV-1's depth-two source
resolvents are sound but incomplete and cubic to discover exhaustively;
LSPV-2's mixed-interface projection is exact only after a verified factorisation
and usually fails the physical length test; LSPV-3's source-relative proof cache
has a correct typed transfer/rent-or-buy invariant but can retain parent-scale
state; LSPV-4 supplies the exact joint exponent gate but no constructor.  None
is silently promoted.

\section{Reproducible finite controls}
The producer exhaustively enumerated $r=1,\ldots,5$ in independent chunks, each
with at most $2^{10}$ ambient points.  It evaluated 1,364 points and 577,251
linear-character terms.  For $r=5$ the maximum nonzero linear bias is exactly
$1/33$, while the quadratic gap is $31/64$.  Twenty symbolic exact rows through
$r=20$ record integer numerators and denominators.  The producer reported 1,854
assertion groups and 57,871 peak traced Python bytes.  A separately written
checker imported no producer code and passed 129 checks.

The required original-source control used the eight proper signed clauses on
one variable triple.  H8 is UNSAT; deleting the clause falsified by $001$ gives
H7 with exactly that one model.  The same DPLL used 7 nodes/180 literal tests on
H8 and 3 nodes/66 literal tests on H7.  For every original signed occurrence the
raw output records $S_j$, $U_j$, and both inversions; all 49 distinct source-pair
SS/SU/US/UU matrices, the coherent $0^3,1^3$ endpoints, two intermediate flips,
per-clause truth counts, and complement identities were independently replayed.
These controls authenticate both labels and guard against confusing local
occurrence targets with global assignments.  They do not connect the bent
distribution to SAT hardness.  Python counts are not native FLMV resources.

\section{Corpus reconciliation and barrier audits}
\paragraph{CORPUS REVIEWED.}
Current root \texttt{main.tex}, AGENTS, CODEX, the pnp-research skill and
available task specifications; REDIRECT1, universal FLMV CONTRACT1,
AFFINE-SMALLBIAS-DESCENT-GATE2, CONSTRUCTION-STRENGTHEN17 and
ARC-LINEAR22-SATURATION1; issue 3; current LSPV-1 through LSPV-4 reports;
LSPV-5 MAIN228--236; the earlier RESTRICTION-SHRINKAGE-BARRIER63,
KRAFT-CERTIFIED-COVER252, ALL-LEAF-OUTPUT-TARIFF264 and hitting-cube audits;
and cited primary literature.

\paragraph{DUPLICATION / CONTRADICTION AUDIT.}
Bent spectra and small-bias definitions are classical.  The project had already
stated that genuinely nonlinear maps escape the affine theorem; this run adds
the exact degree-two countercontrol and proof, rather than repeating the semantic
selector.  The bounded-occurrence cube-cover idea was rejected as duplicate of
the verified Barrier63/Kraft252/Output264 chain and was not republished.

\paragraph{WHAT CHANGED RELATIVE TO ROOT MAIN.}
The synthesis now has an exact firewall at the first nonlinear degree: linear
small bias alone cannot support a quadratic no-go.  The current joint FLMV slack
from LSPV-4 is reconciled without claiming any compiler.

\paragraph{ADOPTED / REJECTED / RETIRED PRIOR COMPONENTS.}
Adopted: the corrected affine counting bound, fixed-seed/source-independence
quantifiers, exact FLMV same-machine gate, source-indexed S/U controls, and
typed proof/cache safeguards.  Rejected: extending affine Fourier fibers to
quadratic fibers without a new pseudorandomness theorem.  Retired: treating a
nonlinear semantic distinguisher as a charged SAT constructor.

\paragraph{STALE OR OVERSTATED CLAIMS FOUND.}
The root header still names the preceding 00:48 snapshot although LSPV-1--4 have
subsequently updated owned regions; whole-document LSPV-5 reconciliation is
therefore due.  No verified report supplies the all-input hereditary child.

\paragraph{Relativization, algebrization, and PQ transfer.}
A SAT oracle constructs the constant semantic selector immediately, so a
universal construction lower bound that relativizes cannot close ARC.  The
quadratic map is exactly an algebraic-access escape from a linear-character
argument.  Schur-square/quadratic-relation language is an exact object analogy
for the degree jump, but no code-based hardness theorem transfers.  Status:
\textbf{FAILED TRANSFER} for arbitrary SAT hardness; \textbf{THEOREM TRANSFER}
only for the elementary Walsh calculation.  No Natural-Proofs conclusion or
cryptographic assumption is used.

\section{Conclusion and next lemma}
\textbf{Proved claim, exact scope:} there are distributions with linear bias
$1/(2^r+1)$ that a fixed quadratic bit distinguishes from uniform with gap
$(1-2^{-r})/2$.  Therefore AFFINE-SMALLBIAS cannot be extended to quadratic maps
from small bias alone.  This neither refutes the affine result nor proves a SAT
runtime bound.

\textbf{Next specific lemma:} for the actual CHMY local easy-source ensemble,
determine and prove the largest degree $d\ge2$ (if any) for which all degree-$d$
output fibers of the proposed transformation are fooled strongly enough after
the complete target-YES union bound, while simultaneously giving a source-
computable both-label SAT transformation and charging its evaluation.  Without
that construction-specific lemma, the nonlinear ARC arrow is OPEN.  P versus NP
remains \textbf{OPEN}.

\begin{thebibliography}{9}
\bibitem{Rothaus} O. S. Rothaus, ``On bent functions,'' \emph{Journal of
Combinatorial Theory, Series A} 20(3), 300--305 (1976),
\url{https://doi.org/10.1016/0097-3165(76)90024-8}.
\bibitem{NaorNaor} J. Naor and M. Naor, ``Small-Bias Probability Spaces:
Efficient Constructions and Applications,'' \emph{SIAM Journal on Computing}
22(4), 838--856 (1993), \url{https://doi.org/10.1137/0222053}.
\bibitem{CHMY} L. Chen, S. Hirahara, R. Meka and R. Williams, ``Lower Bounds for
MCSP and Turing Machines,'' ECCC TR19-022, Theorem 33,
\url{https://eccc.weizmann.ac.il/report/2019/022/}.
\bibitem{FLMV} L. Fortnow, R. Lipton, D. van Melkebeek and A. Viglas,
``Time-Space Lower Bounds for Satisfiability,'' \emph{JACM} 52(6), 835--865
(2005), \url{https://doi.org/10.1145/1101821.1101822}.
\end{thebibliography}
\end{document}
% END_OF_VERIFIED_REPORT::LSPV-5::AFFINE-TO-QUADRATIC-SMALLBIAS-FIREWALL237
